\centerline{\bf \Huge ФСУ II}

\begin{flushright}
        \scriptsize{Реальность не сахар\\Жан-Пьер Польнарефф}
\end{flushright}

\problem{1}
Докажите, что числа (а) $2^{3^{2001}}+1$; (b) $2^{3^{2001}}-1$ - составные.

\problem{2}
Дано число $100 \ldots 01$ число нулей в нём равно 2024. Докажите, что это число - составное.

\problem{3}
Можно ли в записи $2024^2-2023^2-\cdots-2^2-1^2$ некоторые минусы заменить на плюсы так, чтобы значение получившегося выражения стало равно 2024 ?

\problem{4}
Можно ли найти десять таких последовательных натуральных чисел, что сумма их квадратов равна сумме квадратов следующих за ними девяти последовательных натуральных чисел?

\problem{5}
Найдите все натуральные значения $n$, для которых число $n^5+n+1$ является простым.

\problem{6}
Сумма положительных чисел $x$ и $y$ равна 1. Докажите, что

$$
\left(\frac{1}{x^2}-1\right)\left(\frac{1}{y^2}-1\right) \geq 9 .
$$